\chapter{Симуляция и численные эксперименты}\label{ch:simulation}

% TODO (M5/M6): ~5 страниц.
% Структура:
%   6.1 Архитектура симулятора (S3 гибрид)
%       — ODE-модель из §3 для одного пикселя (scipy.integrate.solve_ivp);
%       — surrogate-модель (spline или NN) поверх калибровочного набора;
%       — стек уровней: physical (ODE) → algorithmic (surrogate) → validation (real stand).
%   6.2 Калибровочный pipeline
%       — данные собраны опкодом 0x0E (BENCH_RUN) для baseline B0 и B1;
%       — параметры η, R, q, m из §3.5;
%       — surrogate обучается на расширенном наборе LUT.
%   6.3 Тест-набор (10+ сценариев)
%       — диагностические: BB→WW, чекерборд, текст (LCD font);
%       — реалистичные: 5-6 изображений (графики, диаграммы, фото-эскиз).
%   6.4 Парето-фронт симулированный
%       — наша граница vs точки B0, B1, Kang 2025 (если воспроизводимо), Lin 2024;
%       — рисунки: 2D-проекции (E,G), (E,τ), (G,τ); 3D-облако;
%       — таблица с числами (% улучшения по каждой оси).
%   6.5 Анализ чувствительности
%       — изменение параметров калибровки → смещение Парето;
%       — обоснование робастности выводов.

\section*{Заглушка}

Раздел будет наполнен на этапе~M5/M6 в соответствии с Chapter Plan
(см. \texttt{.ai-factory/RESEARCH.md}).
